\documentclass[11pt,a4paper,parskip=half]{scrartcl} %https://ctan.org/pkg/scrartcl : class script-article
\usepackage[T1]{fontenc}% :allows the user to select font encodings,
\usepackage[ngerman]{babel}%https://ctan.org/pkg/
\usepackage{lmodern}%https://ctan.org/pkg/
\usepackage{microtype}%https://ctan.org/pkg/
\usepackage[margin=23mm,headheight=16pt]{geometry}%https://ctan.org/pkg/
\usepackage{xcolor}%https://ctan.org/pkg/
\usepackage{booktabs}%https://ctan.org/pkg/
\usepackage{tabularx}%https://ctan.org/pkg/
\usepackage{longtable}%https://ctan.org/pkg/
\usepackage{array}%https://ctan.org/pkg/
\usepackage{enumitem}%https://ctan.org/pkg/
\usepackage{csquotes}%https://ctan.org/pkg/
\usepackage{listings}%https://ctan.org/pkg/
\usepackage[most]{tcolorbox}%https://ctan.org/pkg/
\usepackage{amsmath}%https://ctan.org/pkg/
\usepackage{siunitx}%https://ctan.org/pkg/
\usepackage[automark,headsepline]{scrlayer-scrpage}%https://ctan.org/pkg/
\usepackage{hyperref}%https://ctan.org/pkg/
\sisetup{locale=DE}

\definecolor{SFNnavy}{HTML}{18354F}
\definecolor{SFNteal}{HTML}{147F88}
\definecolor{SFNlight}{HTML}{F1F6F8}
\definecolor{SFNcode}{HTML}{F7F9FA}
\hypersetup{colorlinks=true,linkcolor=SFNnavy,urlcolor=SFNteal,pdftitle={LaTeX-Rezeptbuch fuer Forschungsarbeiten}}
\setkomafont{section}{\color{SFNnavy}\Large\bfseries}
\setkomafont{subsection}{\color{SFNnavy}\large\bfseries}
\setkomafont{disposition}{\color{SFNnavy}\bfseries}
\clearpairofpagestyles
\ihead{SFN · LaTeX-Rezeptbuch}
\ohead{\headmark}
\cfoot{\pagemark}
\setlist{itemsep=2pt,topsep=2pt}
\newcolumntype{Y}{>{\raggedright\arraybackslash}X}
\newcommand{\pkg}[1]{\texttt{#1}}
\newcommand{\rezept}[1]{\textsf{#1}}
\lstdefinestyle{latexbook}{
  language=[LaTeX]TeX,
  basicstyle=\ttfamily\footnotesize,
  keywordstyle=\color{SFNteal},
  commentstyle=\color{black!52},
  stringstyle=\color{SFNnavy},
  backgroundcolor=\color{SFNcode},
  frame=single, rulecolor=\color{black!20},
  framesep=7pt, breaklines=true, breakatwhitespace=false,
  columns=fullflexible, keepspaces=true,
  showstringspaces=false,
  numbers=left, numberstyle=\tiny\color{black!45},
  xleftmargin=1.5em, framexleftmargin=1.5em,
  literate=
  {ä}{{\"a}}1 {ö}{{\"o}}1 {ü}{{\"u}}1 {ß}{{\ss}}1
  {Ä}{{\"A}}1 {Ö}{{\"O}}1 {Ü}{{\"U}}1
}
\lstset{style=latexbook}
\newtcolorbox{merkbox}[1][]{
  enhanced,breakable,colback=SFNlight,colframe=SFNteal,
  boxrule=.6pt,arc=1.5mm,left=3mm,right=3mm,
  fonttitle=\bfseries,title={Merke},#1}
\newtcolorbox{aufgabebox}[1][]{
  enhanced,breakable,colback=white,colframe=SFNnavy!70,
  boxrule=.6pt,arc=1.5mm,left=3mm,right=3mm,
  fonttitle=\bfseries,title={Selbst ausprobieren},#1}
\newcommand{\mwe}[2]{%
  \noindent\textbf{MWE-Datei (separat kompilierbar):} \path{#1}\par\vspace{-4pt}
  \lstinputlisting[caption={#2},captionpos=b]{#1}
}

\begin{document}
\begin{titlepage}
  \thispagestyle{empty}
  \vspace*{2cm}
  {\color{SFNteal}\rule{3.5cm}{2.2pt}}\par\vspace{1cm}
  {\Huge\bfseries\color{SFNnavy}LaTeX-Rezeptbuch\par}
  \vspace{.35cm}
  {\LARGE Für naturwissenschaftliche Forschungsarbeiten\par}
  \vspace{1.5cm}
  {\Large Vier Workshoptermine · jeweils 120 Minuten\par}
  \vspace{1cm}
  Vom Titelblatt über Abbildungen, Messwerttabellen und Gleichungen
  bis zu korrekt belegten Quellen und einem vollständigen Dokument.
  \vfill
  \begin{tcolorbox}[colback=SFNlight,colframe=SFNteal,boxrule=.5pt,arc=2mm]
  \textbf{Arbeitsprinzip:} Ein Rezept lesen · die Datei kompilieren · eine
  Änderung ausprobieren · den Baustein in die eigene Forschungsarbeit übernehmen.
  \end{tcolorbox}
  \vspace{1.4cm}
  \textbf{SFN LaTeX-Workshop}\hfill Stand: 28. Oktober 2026
\end{titlepage}

\tableofcontents
\clearpage

\section{So benutzt du dieses Buch}
Dieses Buch ist \emph{kein} Lehrtext, den du von vorne bis hinten abschreiben sollst.
Es ist eine Sammlung von \textbf{Rezepten}: Jedes Rezept löst genau ein typisches
Problem beim Schreiben einer Forschungsarbeit. Der Quelltext steht immer direkt
im Buch und liegt außerdem als eigenständige Datei im Ordner \texttt{rezepte/}.

\begin{merkbox}
\textbf{MWE} bedeutet \emph{Minimal Working Example}: ein möglichst kleines,
aber funktionierendes LaTeX-Dokument mit \texttt{\textbackslash documentclass},
Präambel und \texttt{document}-Umgebung. Es lässt sich unabhängig von den anderen
Beispielen kompilieren, sofern die mitgelieferten Bilder bzw. Daten am angegebenen
Ort liegen. Kopiere für dein Projekt nur die benötigten Pakete in die
\textbf{Präambel} und den eigentlichen Inhalt in den \textbf{Dokumentkörper}.
\end{merkbox}

\textbf{Ordnerstruktur:}
\begin{lstlisting}[language={},numbers=none]
rezeptbuch.tex                 dieses Buch
beispiel-forschungsarbeit.tex komplettes Muster
rezepte/                      13 einzelne MWEs
literatur.bib                 Beispiel-Quellen
bilder/messaufbau.png         frei nutzbares Schema
daten/messwerte.csv           erfundene Messwerte
kapitel/anhang_messprotokoll.tex  ausgelagerter Anhang
README.md                     Startanleitung
\end{lstlisting}

\textbf{Kompilieren:} Öffne den \emph{gesamten} Projektordner in VS Code.
Die im Folgenden angegebenen Terminalbefehle werden im \emph{Projekt-Hauptordner}
ausgeführt, weil Bild-, Daten- und Bib-Dateien über relative Pfade verknüpft sind.
Ein Beispiel ohne Literaturverzeichnis benötigt in der Regel zwei LaTeX-Läufe
für das Inhaltsverzeichnis und die Querverweise. Das Quellenrezept braucht Biber
(siehe Rezept~\ref{subsec:biber}).

\section{Welche Elemente gehören in eine Forschungsarbeit?}
\label{sec:elemente}
Die folgende Übersicht ist zugleich eine \textbf{Packliste}. Nicht jede Arbeit
braucht jedes Element; maßgeblich sind die Vorgaben der betreuenden Schule bzw.
des Wettbewerbs. \textbf{P0} ist unverzichtbare Grundlage, \textbf{P1} häufig
nützlich, \textbf{P2} eine freiwillige Vertiefung.

{\footnotesize
\setlength{\tabcolsep}{4pt}
\setlength{\LTpre}{4pt}\setlength{\LTpost}{8pt}
\begin{longtable}{@{}p{.24\linewidth}p{.30\linewidth}p{.28\linewidth}p{.08\linewidth}@{}}
\toprule
\textbf{Element / Objekt} & \textbf{Befehle / Umgebungen} & \textbf{Paket / Zuständigkeit} & \textbf{Prio}\\
\midrule\endfirsthead
\toprule
\textbf{Element / Objekt} & \textbf{Befehle / Umgebungen} & \textbf{Paket / Zuständigkeit} & \textbf{Prio}\\
\midrule\endhead
Dokumentenklasse, Seitenformat & \texttt{\textbackslash documentclass} & \pkg{scrartcl} (KOMA-Script), bei Bedarf \pkg{geometry} & P0\\
Titelblatt, Metadaten & \texttt{titlepage}, \texttt{\textbackslash centering} & LaTeX/KOMA-Script & P0\\
Zusammenfassung / Abstract & \texttt{abstract} & Dokumentenklasse & P0\\
Kapitel und Gliederung & \texttt{\textbackslash section}, \texttt{\textbackslash subsection} & LaTeX/KOMA-Script & P0\\
Inhaltsverzeichnis & \texttt{\textbackslash tableofcontents} & LaTeX/KOMA-Script & P0\\
Absätze und Hervorhebung & Leerzeile, \texttt{\textbackslash emph} & LaTeX; \pkg{microtype} für Feinsatz & P0\\
Nummerierte / unnummerierte Listen & \texttt{enumerate}, \texttt{itemize} & LaTeX; \pkg{enumitem} zur Anpassung & P0\\
Zitate / Fußnoten & \texttt{\textbackslash enquote}, \texttt{\textbackslash footnote} & \pkg{csquotes}, LaTeX & P1\\
Grafiken / Fotos / Skizzen & \texttt{\textbackslash includegraphics} & \pkg{graphicx} & P0\\
Bildunterschriften & \texttt{figure}, \texttt{\textbackslash caption} & LaTeX; optional \pkg{caption} & P0\\
Tabellen und Beschriftungen & \texttt{table}, \texttt{tabular} & \pkg{booktabs}, \pkg{tabularx} & P0\\
Querverweise & \texttt{\textbackslash label}, \texttt{\textbackslash ref}, \texttt{\textbackslash eqref} & LaTeX; \pkg{hyperref}, optional \pkg{cleveref} & P0\\
Mathematische Formeln & \texttt{equation}, \texttt{align} & \pkg{amsmath}, ggf. \pkg{amssymb} & P0*\\
Messwerte und SI-Einheiten & \texttt{\textbackslash num}, \texttt{\textbackslash qty}, \texttt{\textbackslash si} & \pkg{siunitx} v3 & P0*\\
Quelltext / Programmlistings & \texttt{lstlisting} & \pkg{listings} & P1\\
Diagramme aus Messdaten & \texttt{tikzpicture}, \texttt{axis} & \pkg{pgfplots} & P1\\
Belege / Zitate im Text & \texttt{\textbackslash parencite}, \texttt{\textbackslash textcite} & \pkg{biblatex} und \pkg{biber} & P0\\
Quellenverzeichnis & \texttt{\textbackslash printbibliography} & \pkg{biblatex}, externe \texttt{.bib}-Datei & P0\\
Anhänge / Dateien aufteilen & \texttt{\textbackslash appendix}, \texttt{\textbackslash input} & LaTeX & P1\\
Unterabbildungen & \texttt{subfigure} & \pkg{subcaption} & P2\\
Kopf- und Fußzeilen & \texttt{\textbackslash ihead}, \texttt{\textbackslash cfoot} & \pkg{scrlayer-scrpage} & P2\\
Abbildungs- / Tabellenverzeichnis & \texttt{\textbackslash listoffigures}, \texttt{\textbackslash listoftables} & LaTeX & P2\\
Lange Tabellen über Seiten & \texttt{longtable} & \pkg{longtable} & P2\\
Eigenständigkeitserklärung & \texttt{\textbackslash section*} & LaTeX; nur wenn vorgeschrieben & P1\\
Akronyme / Abkürzungen & \texttt{description}, ggf. Glossarpaket & zunächst LaTeX, später \pkg{glossaries-extra} & P2\\
\bottomrule
\end{longtable}
}
\noindent\textit{* P0 für quantitativ-naturwissenschaftliche Arbeiten; in rein
beschreibenden Arbeiten ggf. nicht nötig.}

\section{Die Reihenfolge in vier Terminen}
\label{sec:termine}
Die Workshopzeit umfasst \textbf{8 Stunden insgesamt}. In jeder Sitzung
entsteht ein neuer Teil \emph{derselben} Muster-Forschungsarbeit. Die MWEs sind
zum Nachschlagen gedacht; im Unterricht werden nicht alle einzeln abgeschrieben.

\begin{tcolorbox}[colback=SFNlight,colframe=SFNteal,title={Termin 1 · Das lesbare Dokument (120 Min.)},fonttitle=\bfseries]
\textbf{0--15} Forschungsarbeit als Dokumentstruktur verstehen ·
\textbf{15--40} Titelseite, Abstract, Grundgerüst (R01) ·
\textbf{40--65} Überschriften, Inhaltsverzeichnis, erste Labels (R02) ·
\textbf{65--85} Absätze, Listen und Zitate (R03) ·
\textbf{85--115} eigene Einleitung und Methodenstruktur erstellen ·
\textbf{115--120} Sicherung und Ausblick.\\
\textbf{Ergebnis:} Titelblatt, Abstract, Inhaltsverzeichnis, vier bis fünf
fachlich sinnvolle Abschnitte und erste Textseiten.
\end{tcolorbox}
\begin{tcolorbox}[colback=SFNlight,colframe=SFNteal,title={Termin 2 · Ergebnisse sichtbar machen (120 Min.)},fonttitle=\bfseries]
\textbf{0--15} Wiederholung und Dateistruktur ·
\textbf{15--55} Abbildungen, Beschriftungen, \texttt{label/ref} (R04) ·
\textbf{55--85} saubere Messwerttabellen (R05) ·
\textbf{85--110} eigene Abbildung und eigene Tabelle einsetzen ·
\textbf{110--120} typische Fehler beheben und Zwischenstand prüfen.\\
\textbf{Ergebnis:} mindestens eine beschriftete Abbildung und eine Tabelle,
beide im Fließtext referenziert. Schnelle Teams: R12.
\end{tcolorbox}
\begin{tcolorbox}[colback=SFNlight,colframe=SFNteal,title={Termin 3 · Wissenschaftliche Inhalte setzen (120 Min.)},fonttitle=\bfseries]
\textbf{0--15} Rückblick: Ergebnisse und Referenzen ·
\textbf{15--45} Formeln und Gleichungsnummern (R06) ·
\textbf{45--65} Zahlen, Messwerte und Einheiten (R07) ·
\textbf{65--85} Programmcodes oder Auswertungsmethode (R08) ·
\textbf{85--110} eigene Auswertung schreiben ·
\textbf{110--120} Literaturdatenbank für Termin 4 anlegen.\\
\textbf{Ergebnis:} nachvollziehbare Methode und Auswertung mit konsistenten
Einheiten; Code nur, wenn fachlich relevant. Optional R10 (CSV-Diagramm).
\end{tcolorbox}
\begin{tcolorbox}[colback=SFNlight,colframe=SFNteal,title={Termin 4 · Quellen, Anhang, Abschluss (120 Min.)},fonttitle=\bfseries]
\textbf{0--15} Was muss belegt werden? ·
\textbf{15--45} \texttt{.bib}, Zitieren und Biber-Lauf (R09) ·
\textbf{45--60} Anhang, Mehrdateienprinzip (R11) ·
\textbf{60--105} vollständige Forschungsarbeit zusammensetzen und prüfen ·
\textbf{105--120} Peer-Review, PDF-Endkontrolle, Abgabe-Check.\\
\textbf{Ergebnis:} vollständiges kompilierbares Dokument vom Titelblatt bis
zum Quellenverzeichnis, bei Bedarf mit Anhang.
\end{tcolorbox}

\begin{merkbox}
\textbf{Didaktische Entscheidung:} Zuerst die \emph{Dokumentstruktur}, dann
\emph{wissenschaftliche Objekte}, erst danach \emph{Feingestaltung}.
Querverweise werden parallel zu Abbildungen und Tabellen eingeführt, nicht als
isoliertes Thema. Quellenangaben erhalten einen eigenen Arbeitsblock, da die
zusätzliche Biber-Verarbeitung Fehlerquellen mitbringt. Für schnellere Gruppen
liegen optionale Rezepte bereit.
\end{merkbox}

\clearpage
\section{Termin 1: Grundgerüst und Text}

\subsection{R01 · Titelblatt, Zusammenfassung und Inhaltsverzeichnis}
\textbf{Du brauchst:} \pkg{scrartcl}, \pkg{babel}, \pkg{microtype}; für klickbare
Links \pkg{hyperref}. \textbf{Warum:} Die Struktur wird später automatisch
verwaltet. Im Titelblatt stehen frei anpassbare Platzhalter.
\mwe{rezepte/01_grundgeruest.tex}{R01: eine komplette, kurze Forschungsarbeit als Startgerüst}
\begin{aufgabebox}
Ändere den Titel und die Forschungsfrage; ergänze unter
\enquote{Methode} zwei Unterabschnitte und erzeuge ein neues PDF.
Prüfe, ob das Inhaltsverzeichnis aktualisiert wird.
\end{aufgabebox}

\subsection{R02 · Überschriften und Querverweise}
\textbf{Du brauchst:} \texttt{\textbackslash section},
\texttt{\textbackslash subsection}, \texttt{\textbackslash label},
\texttt{\textbackslash ref}; \pkg{hyperref} macht Verweise klickbar.
Setze das Label \emph{direkt hinter} die Überschrift.
\mwe{rezepte/02_gliederung.tex}{R02: Gliederung und Abschnittsverweis}
\begin{merkbox}
\textbf{Keine Abschnittsnummern von Hand schreiben.} Bei einer Umstellung der
Gliederung ändern sich Nummerierung und Querverweise automatisch. Das gelingt
in der Regel erst nach einem weiteren LaTeX-Lauf.
\end{merkbox}

\subsection{R03 · Absätze, Zitate, Aufzählungen und Fußnoten}
\textbf{Du brauchst:} \pkg{csquotes} für sprachabhängige Anführungszeichen;
\pkg{enumitem} für saubere nummerierte Listen. Eine Leerzeile beginnt einen
neuen Absatz; \texttt{\textbackslash\textbackslash} ist \emph{kein} Ersatz.
\mwe{rezepte/03_text_und_listen.tex}{R03: verständlicher Fließtext und strukturierte Listen}
\begin{aufgabebox}
Schreibe in drei Sätzen die Hypothese deiner Arbeit.
Füge anschließend eine Materialliste und einen nummerierten Versuchsablauf hinzu.
\end{aufgabebox}

\section{Termin 2: Abbildungen und Tabellen}
\subsection{R04 · Ein Bild mit Quellenbezug und Referenz}
\textbf{Du brauchst:} \pkg{graphicx}; eine \texttt{figure}-Umgebung,
\texttt{\textbackslash caption} und \texttt{\textbackslash label}.
\texttt{[htbp]} bedeutet nur bevorzugte Platzierung, keine feste Position.
Beschreibe im Text, \emph{warum} die Abbildung relevant ist.
\mwe{rezepte/04_abbildung.tex}{R04: Bild einfügen, beschriften und im Text referenzieren}
\begin{merkbox}
Bilder benötigen eine Herkunftsangabe, wenn du sie nicht selbst erstellt hast.
Die \emph{Beschriftung} ist nicht automatisch ein urheberrechtlicher
Nutzungsnachweis. Nutze eigene Bilder oder zulässig lizenzierte Grafiken und
notiere ggf. Lizenz sowie Urheber. Das mitgelieferte Schema ist eigens für
unseren Kurs gezeichnet.
\end{merkbox}

\subsection{R05 · Eine lesbare Messwerttabelle}
\textbf{Du brauchst:} \pkg{booktabs} für klare horizontale Linien,
\pkg{tabularx} für anpassbare Spaltenbreite und \pkg{siunitx} zum
Ausrichten von Zahlen. \texttt{S}-Spalten erwarten Eingaben wie
\texttt{65.0}, geben bei \texttt{locale=DE} aber \num{65.0} aus.
\mwe{rezepte/05_tabelle.tex}{R05: typografisch saubere Messwerttabelle}

\begin{tcolorbox}[colback=white,colframe=black!30,title={So sieht eine Messwerttabelle aus},fonttitle=\bfseries]
\centering
\begin{tabular}{lrr}
\toprule
Messpunkt & Zeit (min) & Temperatur (°C)\\
\midrule
Start & 0 & 65,0\\
Mitte & 10 & 52,4\\
Ende & 20 & 44,1\\
\bottomrule
\end{tabular}
\end{tcolorbox}
\begin{aufgabebox}
Tausche das Beispielbild gegen eine eigene Abbildung aus und übertrage
mindestens drei eigene Messwerte in eine Tabelle. Füge in beiden Fällen
\texttt{caption} und \texttt{label} hinzu und referenziere sie im Text.
\end{aufgabebox}

\section{Termin 3: Formeln, Einheiten und Methoden}
\subsection{R06 · Mathematische Formeln mit Nummern}
\textbf{Du brauchst:} \pkg{amsmath}. Für einzelne Formeln verwende
\texttt{equation}; mehrere zusammengehörige Rechenschritte werden mit
\texttt{align} ausgerichtet. \texttt{\textbackslash eqref} setzt passende Klammern.
\mwe{rezepte/06_formeln.tex}{R06: Mittelwert, Summenzeichen und Formelverweis}
\noindent\textbf{Als Satzbild:}
\begin{equation}\label{eq:demo}
\overline T=\frac{1}{n}\sum_{i=1}^{n} T_i.
\end{equation}

\subsection{R07 · Zahlen, Messwerte und SI-Einheiten}
\textbf{Du brauchst:} \pkg{siunitx} \emph{Version 3}.
Die heute empfohlene Schreibweise ist \texttt{\textbackslash qty}, nicht die
ältere \texttt{\textbackslash SI}. Mit \texttt{locale=DE} werden
Dezimalkommas korrekt gesetzt.
\mwe{rezepte/07_einheiten.tex}{R07: Größen, Zahlen und SI-Einheiten normgerecht setzen}

\subsection{R08 · Programmcode dokumentieren}
\textbf{Du brauchst:} \pkg{listings}. Verwende diesen Baustein nur, wenn
Code zur Reproduzierbarkeit der Arbeit beiträgt. Lange Listings gehören oft in
den Anhang oder ein Repository.
\mwe{rezepte/08_programmcode.tex}{R08: Python-Code mit Zeilennummern darstellen}

\subsection{R10 · Aus einer CSV-Datei ein Diagramm erzeugen (optional)}
\textbf{Du brauchst:} \pkg{pgfplots}. Das Paket liest Zahlen aus einer
Textdatei; bei Änderungen an der Datei entsteht das Diagramm beim nächsten
Kompilieren neu. Für Einsteiger ist auch ein aus Python exportiertes PDF/PNG
plus \pkg{graphicx} völlig ausreichend.
\mwe{rezepte/10_diagramm.tex}{R10: reproduzierbares Diagramm aus fünf CSV-Messwerten}
\begin{aufgabebox}
Füge in die Beispielarbeit eine eigene Gleichung oder einen Rechenschritt ein.
Verwende für jeden numerischen Messwert im Fließtext eine passende Einheit.
\end{aufgabebox}

\section{Termin 4: Quellen, Anhang und vollständige Abgabe}
\subsection{R09 · Quellen sammeln, zitieren und drucken}
\label{subsec:biber}
\textbf{Du brauchst:} \pkg{biblatex} (Zitierstil), das Programm \pkg{biber}
(Bibliografieverarbeitung), \pkg{csquotes} und eine Datei \texttt{literatur.bib}.
Quellenschlüssel wie \texttt{jcgm2008} verweisen auf einen Eintrag in der
Literaturdatenbank. Die Angaben in unserer \texttt{.bib}-Datei sind Beispiele:
Für die eigene Forschungsarbeit müssen sie durch tatsächlich verwendete
Originalquellen ergänzt bzw. ersetzt werden.
\mwe{rezepte/09_quellen.tex}{R09: Literaturzitate und automatisches Quellenverzeichnis}

\textbf{Datei \texttt{literatur.bib} (Auszug):}
\lstinputlisting[language={},firstline=1,lastline=7,caption={Ein BibLaTeX-Datensatz},captionpos=b]{literatur.bib}

\textbf{Die nötigen Verarbeitungsschritte, im Hauptordner:}
\begin{lstlisting}[language={},numbers=none]
pdflatex rezepte/09_quellen.tex
biber 09_quellen
pdflatex rezepte/09_quellen.tex
pdflatex rezepte/09_quellen.tex
\end{lstlisting}
Alternativ übernimmt \texttt{latexmk -pdf rezepte/09\_quellen.tex} die
Abhängigkeitsprüfung und ruft Biber üblicherweise automatisch auf.
\begin{merkbox}
\texttt{biblatex} mit \texttt{backend=biber} verwendet \emph{nicht}
\texttt{bibtex}. Erscheint \enquote{Citation undefined} oder kein Verzeichnis,
prüfe zuerst den Biber-Lauf und die Schreibweise des Zitierschlüssels.
Auch Websites, Bildquellen und Datensätze müssen nach den geltenden
Vorgaben nachvollziehbar nachgewiesen werden.
\end{merkbox}

\subsection{R11 · Einen Anhang anlegen und Dateien auslagern}
\textbf{Du brauchst:} keine weiteren Pakete. \texttt{\textbackslash appendix} beginnt die Anhangsgliederung. Mit \texttt{\textbackslash input} kann
ein Kapitel aus einer anderen Datei eingebunden werden, ohne deren Inhalt
zu kopieren.
\mwe{rezepte/11_anhang_mehrere_dateien.tex}{R11: Anhang und ausgelagertes Messprotokoll}
\textbf{Der Inhalt der eingebundenen Datei:}
\lstinputlisting[caption={kapitel/anhang\_messprotokoll.tex},captionpos=b]{kapitel/anhang_messprotokoll.tex}
\begin{merkbox}
Ein Quellenverzeichnis steht in vielen Forschungsarbeiten \emph{vor} dem Anhang.
Andere Einrichtungen verlangen eine andere Reihenfolge. Richte dich nach den
verbindlichen lokalen Vorgaben; LaTeX kann beides.
\end{merkbox}

\subsection{Das vollständige Musterprojekt}
Öffne \texttt{beispiel-forschungsarbeit.tex}. Dort findest du Titelblatt,
Abstract, Inhaltsverzeichnis, eine Abbildung, eine Tabelle, eine Formel,
Zitate, Quellenverzeichnis und einen eingebundenen Anhang in \emph{einem}
Dokument. Es ist eine \textbf{Demonstration mit erfundenen Messwerten} und darf
nicht als eigener wissenschaftlicher Befund ausgegeben werden.

\textbf{Starten (mit Biber):}
\begin{lstlisting}[language={},numbers=none]
latexmk -pdf beispiel-forschungsarbeit.tex
\end{lstlisting}

\begin{aufgabebox}[title={Abschlussaufgabe: eure eigene Forschungsarbeit}]
Erstellt eine durchgängig kompilierbare Arbeit, die folgende Kriterien erfüllt:
\begin{enumerate}
\item Titelblatt mit Name und Thema, Zusammenfassung und Inhaltsverzeichnis.
\item Fachlich sinnvolle Gliederung: Frage, Methode, Ergebnisse,
      Diskussion und Fazit.
\item Mindestens eine Abbildung und eine Messwerttabelle mit Beschriftung
      sowie Referenzen aus dem Text.
\item Wenn fachlich erforderlich: Gleichung, Rechenweg, Zahlen mit Einheiten
      und/oder reproduzierbarer Programmcode.
\item Mindestens zwei \emph{tatsächlich verwendete} Quellen, im Text zitiert
      und im Verzeichnis vorhanden (bei echten Arbeiten nach Quellenlage).
\item Keine offenen LaTeX-Fehler, ungeklärten Verweise oder unbelegten
      fremden Bilder; fertiges PDF kontrolliert.
\end{enumerate}
\end{aufgabebox}

\section{Optionale Rezepte für schnellere Teams}
\subsection{R12 · Zwei Unterabbildungen nebeneinander}
Nützlich für Vorher-nachher-Vergleiche oder unterschiedliche Mikroskopbilder.
\pkg{subcaption} ermöglicht separat beschriftete Teilbilder.
\mwe{rezepte/12_unterabbildungen.tex}{R12: zwei Teilbilder mit eigenen Labels}

\subsection{R13 · Kopfzeile und Seitenzahlen}
Die KOMA-Script-Erweiterung \pkg{scrlayer-scrpage} ist für unsere
\pkg{scrartcl}-Klasse die passende Lösung. Ein überladenes Layout lenkt von
der Forschung ab; setze diese Formatierung erst zum Schluss ein.
\mwe{rezepte/13_kopf_und_fuss.tex}{R13: zurückhaltende Kopf- und Fußzeilen}

\section{Fehlerdiagnose und letzte Qualitätskontrolle}
\noindent\begin{tabularx}{.97\linewidth}{@{}p{.32\linewidth}Y@{}}
\toprule
\textbf{Symptom} & \textbf{Prüfen / beheben}\\
\midrule
\enquote{File not found} & Dateiname, Schreibweise, Projekt-Hauptordner und relativen Pfad prüfen.\\
\enquote{Undefined control sequence} & Fehlt ein Paket? Ist der Befehl korrekt geschrieben? Ist die Paketversion geeignet?\\
\enquote{Reference undefined / ??} & \texttt{label} und \texttt{ref} müssen identisch sein; zwei LaTeX-Läufe durchführen.\\
\enquote{Citation undefined} & Zitierschlüssel und \texttt{.bib}-Pfad prüfen, Biber ausführen und LaTeX erneut starten.\\
Bild erscheint auf anderer Seite & \texttt{figure} ist ein Gleitobjekt. Nicht alle Bilder mit \texttt{[H]} erzwingen; lieber sinnvoll im Text referenzieren.\\
Tabelle läuft über den Rand & Spalten mit \pkg{tabularx} oder geringerer Inhaltstiefe anpassen, nicht einfach winzig skalieren.\\
Übervolle Box & Fehlermeldung \enquote{Overfull \textbackslash hbox} suchen; zu lange Wörter, URLs oder Tabellen prüfen.\\
\bottomrule
\end{tabularx}

\begin{merkbox}[title={Vor der Abgabe abhaken}]
Die Arbeit ist wissenschaftlich klar gegliedert, alle Bilder und Tabellen werden
im Text erklärt, Messwerte besitzen passende Einheiten, fremde Leistungen sind
belegt, Quellenverzeichnis und Zitate sind konsistent, das PDF enthält keine
leeren Platzhalter oder \texttt{??}. Inhaltliche Qualität ist wichtiger als
ein besonders dekoratives Layout.
\end{merkbox}

\section{Weiterführende Dokumentation}
Die offizielle Dokumentation ist die verlässlichste Referenz für Paketbefehle:
\begin{itemize}
  \item \url{https://ctan.org/pkg/koma-script} – Dokumentenklassen und Layout.
  \item \url{https://ctan.org/pkg/graphicx} – Bildimport.
  \item \url{https://ctan.org/pkg/booktabs} – Tabellensatz.
  \item \url{https://ctan.org/pkg/siunitx} – Zahlen und SI-Einheiten.
  \item \url{https://ctan.org/pkg/amsmath} – Mathematischer Formelsatz.
  \item \url{https://ctan.org/pkg/biblatex} – Quellenverzeichnis und Zitate.
  \item \url{https://ctan.org/pkg/pgfplots} – Diagramme aus Daten.
\end{itemize}

\end{document}
